\chapter{Introduction}
\label{ch:introduction}

\section{The Integration Challenge}
\label{sec:integration-challenge}

\textbf{Systems biology aims to understand living organisms as integrated systems}, not collections of isolated parts. A bacterial cell responding to glucose depletion, for instance, involves:

\begin{itemize}
    \item \textbf{Metabolic changes}: Enzyme kinetics shift as substrate concentrations change (milliseconds to seconds)
    \item \textbf{Gene expression responses}: Transcription factors activate alternative pathway genes (minutes to hours)
    \item \textbf{Regulatory feedback}: Metabolites regulate their own synthesis genes (closing the loop)
\end{itemize}

\textbf{Current computational models address these layers separately}:

\begin{itemize}
    \item \textbf{Metabolic models} (e.g., flux balance analysis, SBML-encoded pathways) capture enzyme kinetics but omit gene regulation
    \item \textbf{Gene regulatory networks} (Boolean logic, differential equations) model transcription factors but ignore metabolite levels
    \item \textbf{Signaling pathway models} represent protein interactions but lack biochemical context
\end{itemize}

\textbf{The result is a fragmented view}: Researchers build separate models for metabolism, regulation, and signaling, then struggle to integrate them. As \citet{Kitano2002} argued, ``to understand biology at the system level, we must examine the structure and dynamics of cellular and organismal function, rather than the characteristics of isolated parts.''

\textbf{The gap}: \textbf{No unified formalism spans metabolic biochemistry and gene regulatory logic in a single compositional framework.} Existing approaches either:

\begin{enumerate}
    \item Model metabolism without regulation (missing biological control)
    \item Model regulation abstractly without biochemical grounding (no mass balance)
    \item Combine formalisms ad-hoc (software integration, not formal semantics)
\end{enumerate}

This thesis addresses the gap by \textbf{extending Petri net formalism} to enable integrated multi-scale biological modeling.

\section{Motivating Example: The Lac Operon System}
\label{sec:lac-operon-motivation}

\subsection{Biological Background}
\label{subsec:lac-operon-biology}

The \textbf{lactose operon} (\emph{lac} operon) in \emph{Escherichia coli} is a classic example of gene regulation coupled to metabolism. When lactose is present but glucose is absent, the cell expresses enzymes to metabolize lactose. \textbf{This behavior requires coordinating three layers}:

\paragraph{Metabolic layer (biochemical reactions):}
\begin{itemize}
    \item \textbf{Glucose metabolism}: \ce{Glucose -> Pyruvate} (glycolysis, 10 enzymatic steps)
    \item \textbf{Lactose metabolism}: \ce{Lactose + H2O -> Glucose + Galactose} ($\beta$-galactosidase enzyme)
    \item \textbf{cAMP production}: When glucose is low, adenylate cyclase produces cAMP from ATP
\end{itemize}

\paragraph{Regulatory layer (transcription factors):}
\begin{itemize}
    \item \textbf{CRP protein} (catabolite repressor protein): Binds cAMP $\to$ cAMP-CRP complex (activator)
    \item \textbf{LacI repressor}: Blocks \emph{lac} promoter unless allolactose (lactose metabolite) binds
    \item \textbf{Lac promoter}: Requires \textbf{both} cAMP-CRP binding \textbf{and} LacI derepression for transcription
\end{itemize}

\paragraph{Integration (metabolite-gene coupling):}
\begin{itemize}
    \item \textbf{Metabolite $\to$ Transcription}: Low glucose $\to$ High cAMP $\to$ cAMP-CRP $\to$ \emph{lacZ} transcription
    \item \textbf{Product feedback}: Lactose $\to$ Allolactose $\to$ LacI derepression $\to$ More $\beta$-galactosidase
\end{itemize}

\subsection{Why This Cannot Be Modeled with Existing Formalisms}
\label{subsec:lac-operon-limitations}

\paragraph{Classical Petri nets} model biochemical reactions (\ce{glucose -> pyruvate}) but lack:

\begin{itemize}
    \item \textbf{Regulatory arcs}: Cannot represent ``cAMP-CRP activates transcription'' as an arc (would require code)
    \item \textbf{Threshold logic}: Cannot encode ``LacI blocks promoter when $M(\text{allolactose}) < 0.1$ mM'' in topology
    \item \textbf{Heterogeneous dynamics}: Glycolysis is continuous (ODE), gene expression is stochastic bursts (Gillespie) $\to$ cannot mix in one model
\end{itemize}

\paragraph{SBML} (Systems Biology Markup Language) models metabolic reactions but:

\begin{itemize}
    \item \textbf{Gene regulation requires separate models}: SBML-qual (qualitative) extension for Boolean logic, disconnected from kinetics
    \item \textbf{No compositional integration}: Metabolic SBML + regulatory SBML-qual = two files, manually synchronized
\end{itemize}

\paragraph{Boolean gene regulatory networks} model \emph{lacI} $\to$ \emph{lacZ} logic but:

\begin{itemize}
    \item \textbf{No metabolite states}: Cannot represent cAMP concentration (continuous variable)
    \item \textbf{No stoichiometry}: Cannot track lactose consumption, glucose production
    \item \textbf{No enzyme kinetics}: $\beta$-galactosidase activity is binary (on/off), not Michaelis-Menten
\end{itemize}

\paragraph{ODE systems} can model both layers by writing coupled differential equations:

\begin{equation}
\begin{aligned}
\frac{d[\text{Glucose}]}{dt} &= -v_{\text{glycolysis}} \\
\frac{d[\text{cAMP}]}{dt} &= k_{\text{synth}} \cdot \left(1 - \frac{[\text{Glucose}]}{K}\right) - k_{\text{deg}} \cdot [\text{cAMP}] \\
\frac{d[\text{lacZ mRNA}]}{dt} &= k_{\text{tx}} \cdot [\text{cAMP-CRP}] \cdot (1 - [\text{LacI}]) - k_{\text{deg}} \cdot [\text{lacZ mRNA}]
\end{aligned}
\end{equation}

\textbf{But}:

\begin{itemize}
    \item \textbf{Not compositional}: Adding a new metabolite or gene requires rewriting equations (no modularity)
    \item \textbf{No visual topology}: Regulatory connections hidden in equation code
    \item \textbf{No formal analysis}: Cannot use Petri net structural analysis (P-invariants, reachability)
\end{itemize}

\subsection{What Is Needed}
\label{subsec:lac-operon-requirements}

\textbf{A unified formalism must support}:

\begin{enumerate}
    \item \textbf{Metabolic reactions} with stoichiometry (\ce{Glucose + ATP -> G6P + ADP})
    \item \textbf{Gene expression} with regulatory logic (\emph{lacZ} transcription requires cAMP-CRP and $\neg$LacI)
    \item \textbf{Metabolite-gene coupling} (cAMP concentration $\to$ transcription factor activity)
    \item \textbf{Visual topology} (arcs encode regulation, not hidden in code)
    \item \textbf{Heterogeneous dynamics} (continuous glycolysis + stochastic transcription bursts)
    \item \textbf{Mass balance} (atoms conserved: \ce{C6H12O6 -> 2 C3H4O3})
\end{enumerate}

\textbf{This thesis presents such a formalism}: Extended Biological Petri Nets.

\section{Research Questions}
\label{sec:research-questions}

This thesis addresses six fundamental questions about integrated biological modeling:

\subsection{RQ1: Cooperative Parallelism (Weak Independence)}
\label{subsec:rq1}

\textbf{Can Petri nets support parallel execution when transitions share places?}

\begin{itemize}
    \item \textbf{Classical Petri net theory}: Transitions are \textbf{strongly independent} if they share no places (disjoint neighborhoods) $\to$ can fire simultaneously
    \item \textbf{Biological reality}: Enzymes catalyze multiple reactions (shared catalyst place), pathways converge (shared product place)
    \item \textbf{Example}: Hexokinase catalyzes glucose phosphorylation \textit{and} fructose phosphorylation $\to$ shares enzyme place
    \item \textbf{Question}: Can we define \textbf{weak independence} allowing shared catalysts and products while preserving correctness?
\end{itemize}

\textbf{Hypothesis}: Transitions with \textbf{disjoint input places} (no resource conflicts) but shared outputs or catalysts are \textbf{weakly independent} $\to$ can execute in parallel without violating reachability.

\subsection{RQ2: Heterogeneous Transition Types}
\label{subsec:rq2}

\textbf{Can continuous, stochastic, timed, and burst transitions coexist in a single model with consistent semantics?}

\begin{itemize}
    \item \textbf{Biological systems exhibit multiple temporal scales}:
    \begin{itemize}
        \item \textbf{Continuous}: Enzyme kinetics (Michaelis-Menten), equilibria (ms-s)
        \item \textbf{Stochastic}: Gene expression bursts, rare molecular events (low copy numbers)
        \item \textbf{Timed}: Cell cycle checkpoints, scheduled events
        \item \textbf{Burst}: Transcriptional pulsing (chromatin remodeling)
    \end{itemize}
    \item \textbf{Example}: Energy sensing motif requires continuous glycolysis + stochastic gene transcription bursts
    \item \textbf{Question}: Can these four transition types fire in a unified hybrid simulation without semantic conflicts?
\end{itemize}

\textbf{Hypothesis}: A \textbf{hybrid scheduler} coordinating four specialized engines (ODE, Gillespie, event queue, burst sampler) can simulate heterogeneous transitions with provable correctness.

\subsection{RQ3: Arc-Level Regulation}
\label{subsec:rq3}

\textbf{Can regulatory logic (thresholds, Hill equations, inhibition) be encoded directly on arcs?}

\begin{itemize}
    \item \textbf{Classical Petri nets}: Arcs have weights (stoichiometry) but no regulatory semantics
    \item \textbf{Biological regulation}: ATP inhibits phosphofructokinase when $M(\text{ATP}) \geq 5.0$ mM (allosteric feedback)
    \item \textbf{Current approach}: Implement inhibition in transition code (hidden from topology)
    \item \textbf{Question}: Can arcs carry \textbf{threshold formulas} and \textbf{Hill equations} making regulation visible in network structure?
\end{itemize}

\textbf{Hypothesis}: Extending arcs with \textbf{type} (normal/test/inhibitor) and \textbf{threshold functions} embeds regulatory logic in topology, enabling visual reasoning and formal analysis.

\subsection{RQ4: Atomic Conservation}
\label{subsec:rq4}

\textbf{Can biochemical formulas replace abstract tokens to enable elemental balance analysis?}

\begin{itemize}
    \item \textbf{Classical Petri nets}: Places hold tokens (abstract, countable units)
    \item \textbf{Biochemistry}: Molecules have elemental composition (\ce{C6H12O6} = glucose)
    \item \textbf{Mass balance}: Reactions must conserve atoms (6 carbons in $\to$ 6 carbons out)
    \item \textbf{Question}: Can place names be biochemical formulas enabling automatic stoichiometry validation?
\end{itemize}

\textbf{Hypothesis}: Associating places with \textbf{formulas} (element$\to$count dictionaries) enables \textbf{elemental balance verification} at the Petri net level, detecting modeling errors (unbalanced reactions).

\subsection{RQ5: Biological Validity}
\label{subsec:rq5}

\textbf{Does the extended formalism preserve biological correctness principles?}

\begin{itemize}
    \item \textbf{Enzyme conservation}: Catalysts are not consumed (test arcs must preserve marking)
    \item \textbf{Superposition}: Multiple reactions produce the same metabolite (additive, not conflicting)
    \item \textbf{Stoichiometric precision}: Integer coefficients, no fractional molecules in discrete models
    \item \textbf{Regulatory consistency}: Inhibitor arcs block transitions consistently (threshold always enforced)
    \item \textbf{Question}: Can we prove the formalism respects these biological constraints?
\end{itemize}

\textbf{Hypothesis}: \textbf{Well-formedness constraints} (formal rules on network structure) enforce biological validity mechanically, ruling out invalid models.

\subsection{RQ6: Practical Applicability}
\label{subsec:rq6}

\textbf{Can real biological systems be successfully modeled with the extended formalism?}

\begin{itemize}
    \item \textbf{Validation requirement}: Theoretical formalism is only useful if practitioners can model actual systems
    \item \textbf{Examples}: Glycolysis (10 reactions, 3 regulatory checkpoints), TCA cycle (8 reactions, cyclic), cellular respiration (32 reactions, spanning glycolysis + TCA + oxidative phosphorylation)
    \item \textbf{Question}: Do the four innovations (weak independence, heterogeneous types, arc regulation, formulas) enable modeling complex real pathways?
\end{itemize}

\textbf{Hypothesis}: A \textbf{progressive example series} (16 models from simple ATP hydrolysis to complete cellular respiration) demonstrates practical feasibility and identifies formalism limitations.

\section{Thesis Contributions}
\label{sec:contributions}

This thesis makes \textbf{theoretical, validation, and implementation contributions} to systems biology modeling.

\subsection{Theoretical Contributions (Part II: Core Theory)}
\label{subsec:theoretical-contributions}

\paragraph{Chapter \ref{ch:extended-biopn}: Extended Biological Petri Net Formalism}

\begin{itemize}
    \item \textbf{12-tuple definition}: $(P, T, F, W, M_0, K, \Phi, \Sigma, \Theta, \Delta, \tau, \rho)$
    \begin{itemize}
        \item Classical components $(P, T, F, W, M_0, K)$: Places, transitions, arcs, weights, initial marking, capacities
        \item \textbf{Extensions}:
        \begin{itemize}
            \item $\Phi$: Rate functions (Michaelis-Menten, Hill, mass action, custom)
            \item $\Sigma$: Test arcs (non-consumptive, for catalysis)
            \item $\Theta$: Inhibitor arcs (threshold-based blocking)
            \item $\Delta$: Threshold functions (constant, dynamic, Hill equation)
            \item $\tau$: Transition types (Continuous, Stochastic, Timed, Burst)
            \item $\rho$: Biochemical formulas (element$\to$count dictionaries)
        \end{itemize}
    \end{itemize}
    \item \textbf{Arc semantics}: Firing rules for normal/test/inhibitor arcs across all transition types
    \item \textbf{Well-formedness constraints}: 8 formal rules enforcing biological validity (enzyme conservation, stoichiometric consistency)
\end{itemize}

\paragraph{Chapter \ref{ch:weak-independence}: Weak Independence Theory}

\begin{itemize}
    \item \textbf{Formal definition}: Transitions $t_1, t_2$ are \textbf{weakly independent} if $(\bullet t_1 \cap \bullet t_2) = \emptyset$ (disjoint input places)
    \begin{itemize}
        \item Allows shared outputs (convergent reactions $\to$ superposition)
        \item Allows shared catalysts via test arcs (enzyme serves multiple reactions)
    \end{itemize}
    \item \textbf{Dependency classification algorithm}: $O(|T|^2 \cdot |P|)$ algorithm assigning CONFLICT, COUPLING, or INDEPENDENT to all transition pairs
    \item \textbf{Reachability preservation theorem}: Weakly independent transitions can fire concurrently without affecting reachability set
    \item \textbf{Biological cooperativity}: Formalization of how metabolic pathways cooperate (shared intermediates, shared enzymes)
\end{itemize}

\paragraph{Chapter \ref{ch:biochemical-formulas}: Biochemical Formula Tracking}

\begin{itemize}
    \item \textbf{Formula representation}: Hill notation (\ce{C6H12O6}) $\to$ element-count dictionary \{``C'': 6, ``H'': 12, ``O'': 6\}
    \item \textbf{Elemental balance verification}: Algorithm checking $\sum(\text{atoms}_{\text{in}}) = \sum(\text{atoms}_{\text{out}})$ for all transitions
    \item \textbf{Elemental balance matrix} $S_e$: (elements $\times$ transitions), each entry = net atom change
    \item \textbf{Cofactor suggestion}: Algorithm proposing missing cofactors (\ce{H2O}, \ce{H+}, \ce{Pi}) when reactions are unbalanced
    \item \textbf{Database integration}: KEGG/ChEBI formula retrieval for automatic model enrichment
\end{itemize}

\subsection{Validation Contributions (Part III: Empirical Validation)}
\label{subsec:validation-contributions}

\paragraph{Chapter \ref{ch:validation}: Validation Through Examples}

\begin{itemize}
    \item \textbf{16-example progressive series}: From simple (ATP hydrolysis) to complex (complete cellular respiration)
    \item \textbf{Phase structure}:
    \begin{itemize}
        \item Phase 1 (Examples 01-03): Foundation (basic reactions, reversibility, catalysis)
        \item Phase 2 (Examples 04-06): Regulation (inhibitor arcs, competitive inhibition, feedback)
        \item Phase 3 (Examples 07-08): Integration (multi-step pathways, heterogeneous types)
        \item Phase 4 (Examples 09-13): Complete pathways (glycolysis, TCA, respiration)
        \item Phase 5 (Examples 14-16): Advanced (branching, resource competition, dynamic thresholds)
    \end{itemize}
    \item \textbf{Key proof}: Example 08 (Energy Sensing Motif) demonstrates \textbf{all four innovations simultaneously}:
    \begin{itemize}
        \item Weak independence: PFK and PK have disjoint inputs
        \item Heterogeneous types: Continuous (PFK, PK, ATPase) + Stochastic burst (Gene\_PFK)
        \item Arc regulation: ATP inhibitor arcs on PFK and PK (thresholds 2.5 mM, 2.0 mM)
        \item Atomic conservation: All reactions elementally balanced (verified)
    \end{itemize}
    \item \textbf{Quantitative validation}: Parameters from BRENDA database, simulation results match literature
    \item \textbf{Scalability analysis}: Parallel execution achieves 2-4$\times$ speedup on typical networks (8 cores)
\end{itemize}

\subsection{Implementation Contributions (Part IV: Supporting Tools)}
\label{subsec:implementation-contributions}

\paragraph{Chapter \ref{ch:architecture}: SHYpn System Architecture}
\begin{itemize}
    \item Three-tier architecture (Presentation/Business Logic/Data layers)
    \item Model representation: Faithful 12-tuple implementation in Python
    \item UI components: Network canvas (Cairo rendering), property panels (GTK4)
    \item Persistence: Native JSON format + SBML import + GraphML export
\end{itemize}

\paragraph{Chapter \ref{ch:kegg}: KEGG Compound Integration}
\begin{itemize}
    \item REST API wrapper for KEGG database (18,000+ compounds, 11,000+ reactions)
    \item Automatic formula retrieval: User enters KEGG ID $\to$ formula auto-filled
    \item Reaction import: KEGG reaction ID $\to$ creates places, transition, arcs automatically
    \item Pathway bulk import: Entire glycolysis pathway imported in $<5$ seconds
    \item Cofactor suggestion: Algorithm proposes missing \ce{H2O}, \ce{H+}, \ce{Pi} based on elemental imbalance
    \item Local caching: SQLite database $\to$ 60$\times$ speedup on repeated queries
\end{itemize}

\paragraph{Chapter \ref{ch:brenda}: Intelligent Parameter Inference from BRENDA}
\begin{itemize}
    \item SOAP API integration with BRENDA (2.7 million kinetic parameters)
    \item Statistical aggregation: Median $K_m$ + 95\% confidence interval (robust to outliers)
    \item Context-aware heuristics:
    \begin{itemize}
        \item Organism priority (\emph{Saccharomyces cerevisiae} $>$ human $>$ all)
        \item Substrate fuzzy matching (``glucose'' matches ``D-glucose'', ``$\alpha$-D-glucose'')
        \item Quality filtering (citation presence, experimental conditions, outlier detection)
    \end{itemize}
    \item Local database: Progressive accumulation $\to$ 200$\times$ speedup on cached data
    \item Batch enrichment: Entire glycolysis pathway (10 enzymes) parameterized in $<10$ seconds
\end{itemize}

\paragraph{Chapter \ref{ch:simulation}: Hybrid Simulation Engine}
\begin{itemize}
    \item Four-engine architecture:
    \begin{itemize}
        \item Continuous: SciPy RK45 ODE solver (adaptive time-stepping)
        \item Stochastic: Gillespie SSA (exact algorithm)
        \item Timed: Priority queue for scheduled events
        \item Burst: Exponential burst frequency + geometric burst size
    \end{itemize}
    \item Hybrid scheduler: Event-driven coordination (asks all engines ``when is your next event?'')
    \item Parallel execution: Weak independence-based partitioning $\to$ 2-4$\times$ speedup (8 cores)
    \item Validation: Correctness proofs (reachability preservation), benchmark suite (16 examples)
\end{itemize}

\subsection{Impact Summary}
\label{subsec:impact-summary}

\textbf{This thesis enables, for the first time}:

\begin{enumerate}
    \item \textbf{Parallel simulation exploiting biological cooperativity}: Weak independence theory allows 2-4$\times$ speedup on multi-core systems by executing non-conflicting reactions simultaneously
    \item \textbf{Multi-scale integrated models}: Single model spans continuous enzyme kinetics (glycolysis) and stochastic gene expression bursts (\emph{lacZ} transcription)
    \item \textbf{Topology-embedded regulation}: Inhibitor arcs with threshold formulas make regulatory logic visible in network structure (no hidden code)
    \item \textbf{Atomic-level mass balance}: Biochemical formula tracking enables elemental conservation verification, detecting modeling errors automatically
\end{enumerate}

\textbf{Broader significance}:

\begin{itemize}
    \item \textbf{Computational systems biology}: Provides formal foundation for integrated modeling (metabolomics + transcriptomics)
    \item \textbf{Synthetic biology}: Enables design and simulation of engineered pathways with regulatory feedback
    \item \textbf{Drug discovery}: Models metabolic perturbations (enzyme inhibition) coupled to gene expression responses
    \item \textbf{Education}: Visual formalism makes biochemical regulation accessible to students (arc types visible)
\end{itemize}

\section{Research Methodology}
\label{sec:methodology}

\subsection{Formalism Development}
\label{subsec:formalism-development}

\textbf{Iterative refinement approach}:

\begin{enumerate}
    \item \textbf{Requirements analysis}: Survey biological systems needing integrated modeling (lac operon, glycolysis regulation, etc.)
    \item \textbf{Formalism design}: Extend Petri net 5-tuple to 12-tuple, adding $\Phi, \Sigma, \Theta, \Delta, \tau, \rho$ components
    \item \textbf{Formal specification}: Define enabling conditions, firing rules, semantics for each transition type
    \item \textbf{Theoretical analysis}: Prove weak independence reachability preservation theorem
    \item \textbf{Validation}: Test formalism on 16 progressive examples (simple $\to$ complex)
\end{enumerate}

\textbf{Design criteria}:

\begin{itemize}
    \item \textbf{Minimality}: Add only necessary components (no feature bloat)
    \item \textbf{Compositionality}: Models combine without side effects
    \item \textbf{Biological fidelity}: Respects fundamental biological principles (enzyme conservation, mass balance)
    \item \textbf{Formal analyzability}: Enables structural analysis (P-invariants, reachability, weak independence classification)
\end{itemize}

\subsection{Validation Strategy}
\label{subsec:validation-strategy}

\textbf{Progressive example series} (16 models):

\begin{itemize}
    \item \textbf{Phase 1}: Establish baseline (simple reactions, reversibility, catalysis)
    \item \textbf{Phase 2}: Add regulation (inhibitor arcs, feedback loops)
    \item \textbf{Phase 3}: Demonstrate integration (continuous + stochastic, weak independence)
    \item \textbf{Phase 4}: Scale complexity (complete pathways: glycolysis, TCA, respiration)
    \item \textbf{Phase 5}: Stress-test limits (branching, competition, dynamic thresholds)
\end{itemize}

\textbf{Validation criteria per example}:

\begin{enumerate}
    \item \textbf{Structural correctness}: Satisfies well-formedness constraints (C1-C8)
    \item \textbf{Elemental balance}: All transitions conserve atoms (verified automatically)
    \item \textbf{Parameter realism}: $K_m$, $V_{\max}$ from BRENDA database (not arbitrary)
    \item \textbf{Simulation convergence}: Reaches steady-state or expected dynamics
    \item \textbf{Literature consistency}: Results match published experimental data
\end{enumerate}

\textbf{Key proof example}: Example 08 (Energy Sensing Motif) is the \textbf{minimal model demonstrating all four innovations}. If this example works correctly (which it does, verified), it proves the formalism is viable.

\subsection{Implementation Validation}
\label{subsec:implementation-validation}

\textbf{SHYpn platform serves as proof-of-concept}:

\begin{itemize}
    \item Implements all 12 formalism components
    \item Executes all 16 validation examples successfully
    \item Achieves 2-4$\times$ speedup via parallel execution (8 cores)
    \item Integrates KEGG (formulas) and BRENDA (parameters)
\end{itemize}

\textbf{Implementation is secondary contribution}: Thesis proves formalism is sound; tool proves formalism is executable.

\section{Thesis Scope and Limitations}
\label{sec:scope-limitations}

\subsection{In Scope}
\label{subsec:in-scope}

\textbf{What this thesis addresses}:

\begin{itemize}
    \item \textbf{Formalism}: Extending Petri nets for integrated biological modeling
    \item \textbf{Theory}: Weak independence, heterogeneous transitions, arc regulation, formula tracking
    \item \textbf{Validation}: 16 biochemical examples demonstrating practical applicability
    \item \textbf{Implementation}: SHYpn platform as proof-of-concept
\end{itemize}

\textbf{Biological domains covered}:

\begin{itemize}
    \item Central metabolism (glycolysis, TCA cycle, oxidative phosphorylation)
    \item Gene regulation (transcription, mRNA dynamics, protein production)
    \item Simple feedback loops (product inhibition, feed-forward loops)
\end{itemize}

\subsection{Out of Scope}
\label{subsec:out-of-scope}

\textbf{What this thesis does not address}:

\begin{itemize}
    \item \textbf{Spatial dynamics}: Petri nets model well-mixed compartments (no diffusion, cell geometry)
    \begin{itemize}
        \item \textbf{Limitation}: Cannot model spatial gradients (e.g., morphogen gradients in development)
        \item \textbf{Future work}: Extend to spatially distributed Petri nets
    \end{itemize}
    
    \item \textbf{Multi-compartment systems}: Current formalism is single-compartment
    \begin{itemize}
        \item \textbf{Limitation}: Cannot model cytoplasm-mitochondria interactions with transport
        \item \textbf{Future work}: Add compartment types, transport arcs
    \end{itemize}
    
    \item \textbf{Protein-protein interactions}: Focus is metabolites + genes
    \begin{itemize}
        \item \textbf{Limitation}: Signaling cascades (MAPK, JAK-STAT) not fully explored
        \item \textbf{Future work}: Extend to protein complexes, post-translational modifications
    \end{itemize}
    
    \item \textbf{Evolutionary dynamics}: No mutation, selection, population genetics
    \begin{itemize}
        \item \textbf{Out of scope}: Thesis is mechanistic modeling, not evolutionary
    \end{itemize}
    
    \item \textbf{Whole-cell models}: Validated on pathways (10-50 transitions), not entire genomes
    \begin{itemize}
        \item \textbf{Limitation}: Scalability to 1000+ transition networks unproven
        \item \textbf{Future work}: Hierarchical modeling, abstraction techniques
    \end{itemize}
\end{itemize}

\subsection{Assumptions}
\label{subsec:assumptions}

\textbf{Modeling assumptions}:

\begin{enumerate}
    \item \textbf{Well-mixed compartments}: No spatial heterogeneity (stirred reactor assumption)
    \item \textbf{Deterministic enzyme concentrations}: $[E]_0$ is fixed (no enzyme degradation modeled)
    \item \textbf{Steady-state approximation for complexes}: Enzyme-substrate complex formation is fast
    \item \textbf{Ideal solution}: No crowding effects, activity coefficients = 1
    \item \textbf{Constant pH/temperature}: Environmental parameters fixed
\end{enumerate}

\textbf{These assumptions are standard in systems biology} and do not invalidate the formalism's contributions.

\section{Thesis Organization}
\label{sec:thesis-organization}

\textbf{Part I: Introduction and Foundations} (Chapters \ref{ch:introduction}--\ref{ch:integration-challenge})

\begin{itemize}
    \item \textbf{Chapter \ref{ch:introduction}}: Motivation, research questions, contributions (this chapter)
    \item \textbf{Chapter \ref{ch:background}}: Background (Petri nets, biological modeling, related work)
    \item \textbf{Chapter \ref{ch:integration-challenge}}: Integration challenge (lac operon deep dive, requirements analysis)
\end{itemize}

\textbf{Part II: Core Theory -- The Extended Formalism} (Chapters \ref{ch:extended-biopn}--\ref{ch:biochemical-formulas})

\begin{itemize}
    \item \textbf{Chapter \ref{ch:extended-biopn}}: Extended Bio-PN definition (12-tuple, arc semantics, well-formedness)
    \item \textbf{Chapter \ref{ch:weak-independence}}: Weak independence theory (definition, algorithm, reachability theorem)
    \item \textbf{Chapter \ref{ch:biochemical-formulas}}: Biochemical formula tracking (elemental balance, cofactor suggestion)
\end{itemize}

\textbf{Part III: Validation Through Examples} (Chapter \ref{ch:validation})

\begin{itemize}
    \item \textbf{Chapter \ref{ch:validation}}: 16-example progressive series (Phase 1-5, quantitative validation)
\end{itemize}

\textbf{Part IV: Implementation -- Supporting Tools} (Chapters \ref{ch:architecture}--\ref{ch:simulation})

\begin{itemize}
    \item \textbf{Chapter \ref{ch:architecture}}: SHYpn system architecture (three-tier design, model representation)
    \item \textbf{Chapter \ref{ch:kegg}}: KEGG integration (formula retrieval, reaction import, caching)
    \item \textbf{Chapter \ref{ch:brenda}}: Parameter inference (BRENDA API, statistical aggregation, heuristics)
    \item \textbf{Chapter \ref{ch:simulation}}: Hybrid simulation engine (four engines, scheduler, parallel execution)
\end{itemize}

\textbf{Part V: Evaluation} (Chapters \ref{ch:case-studies}--\ref{ch:performance})

\begin{itemize}
    \item \textbf{Chapter \ref{ch:case-studies}}: Case studies (glycolysis regulation, TCA cycle, cellular respiration)
    \item \textbf{Chapter \ref{ch:performance}}: Performance evaluation (scalability, speedup analysis, comparison)
\end{itemize}

\textbf{Part VI: Synthesis} (Chapters \ref{ch:discussion}--\ref{ch:conclusion})

\begin{itemize}
    \item \textbf{Chapter \ref{ch:discussion}}: Discussion (contributions, limitations, biological validity)
    \item \textbf{Chapter \ref{ch:conclusion}}: Conclusion and future work (research impact, open problems)
\end{itemize}

\textbf{Estimated length}: 180-200 pages (10-12 pages per chapter average)

\section{Reading Guide}
\label{sec:reading-guide}

\textbf{For readers interested in}:

\paragraph{Theoretical foundations:}
\begin{itemize}
    \item Read Chapters \ref{ch:extended-biopn}--\ref{ch:biochemical-formulas} (formalism definition, weak independence, formula tracking)
    \item Skip implementation details (Chapters \ref{ch:architecture}--\ref{ch:simulation})
\end{itemize}

\paragraph{Practical modeling:}
\begin{itemize}
    \item Read Chapter \ref{ch:validation} (validation examples)
    \item Read Chapters \ref{ch:kegg}--\ref{ch:brenda} (KEGG/BRENDA integration for parameter inference)
    \item Skim theory (Chapters \ref{ch:extended-biopn}--\ref{ch:biochemical-formulas}) for terminology
\end{itemize}

\paragraph{Implementation and tools:}
\begin{itemize}
    \item Read Chapters \ref{ch:architecture}--\ref{ch:simulation} (architecture, simulation engine, parallel execution)
    \item Chapter \ref{ch:validation} for example suite
\end{itemize}

\paragraph{Systems biology applications:}
\begin{itemize}
    \item Read Chapter \ref{ch:integration-challenge} (integration challenge, lac operon)
    \item Read Chapter \ref{ch:validation} (validation examples)
    \item Read Chapter \ref{ch:case-studies} (case studies)
\end{itemize}

\paragraph{Computer science formalism:}
\begin{itemize}
    \item Read Chapters \ref{ch:extended-biopn}--\ref{ch:weak-independence} (12-tuple definition, weak independence algorithm)
    \item Chapter \ref{ch:simulation} (hybrid simulation semantics)
\end{itemize}

\section{Summary}
\label{sec:ch1-summary}

\textbf{This chapter established the need for integrated biological modeling}, motivated by the lac operon system spanning metabolism and gene regulation. \textbf{Six research questions} guide the thesis:

\begin{enumerate}
    \item Can Petri nets support parallel execution with shared places? (\textbf{RQ1: Weak independence})
    \item Can continuous, stochastic, timed, and burst transitions coexist? (\textbf{RQ2: Heterogeneity})
    \item Can regulatory logic be encoded on arcs? (\textbf{RQ3: Arc regulation})
    \item Can biochemical formulas enable atomic conservation? (\textbf{RQ4: Formula tracking})
    \item Does the formalism preserve biological correctness? (\textbf{RQ5: Validity})
    \item Can real systems be modeled? (\textbf{RQ6: Applicability})
\end{enumerate}

\textbf{Contributions span theory} (Extended Bio-PN formalism, weak independence theory), \textbf{validation} (16-example progressive series), and \textbf{implementation} (SHYpn platform with KEGG/BRENDA integration).

\textbf{Chapter \ref{ch:background} surveys related work}, establishing context for this thesis's innovations. \textbf{Chapter \ref{ch:integration-challenge} presents the integration challenge} in depth, deriving formal requirements for unified modeling.
